\subsection{The same candidates extend to the full spherical action}

\begin{statement}{D23.P30}{Proved}{Full four-generator spherical covariance leaves all eighteen pairs}
\end{statement}
For every one of D23.P26's eighteen pairs, there is a
unique permutation $V_5$ of the degree-$6720$ orbit
whose restriction to the original sector is the known
$V$, and which satisfies the four literal covariance
equations
\[
\begin{split}
 V_5\sigma_1V_5^{-1}&=\sigma_1,\\
 V_5\sigma_2V_5^{-1}&=F^{-1}\sigma_2F,\\
 V_5\sigma_3V_5^{-1}&=R\sigma_3R^{-1},\\
 V_5\sigma_4V_5^{-1}&=\sigma_4,
\end{split}
\]
where $F$ is the complete $f(X_1,X_2)$ value of P29
and $R=r_5^{-1}Fr_5$ is the complete $f(X_3,X_4)$
value. Every such $V_5$ lies in the independent product
$S^{60}$ proved in D23.F39, and hence in the derived
pure finite image. This supplies no further exclusion
of the eighteen pairs.

\begin{proof}
The rotation marking $r_5X_ir_5^{-1}=X_{i+3}$ shows
that conjugation by $r_5^{-1}$ sends $(X_1,X_2)$ to
$(X_3,X_4)$. Thus the displayed expression for $R$
has the exact intended evaluation marking.
The usual Drinfeld formula gives the first three
displayed images, as in P18. For the last generator
its two conjugator arguments are
\[
 y_4=\sigma_3\sigma_2\sigma_1^2\sigma_2\sigma_3,
                    \qquad \sigma_4^2.
\]
The spherical relation gives $y_4=\sigma_4^{-2}$,
also checked literally in the certificate. A commutator
word evaluated at these commuting elements is one,
so the fourth image is $\sigma_4$ at multiplier one.

Let $g_1,\ldots,g_4$ be the old generators and
$g'_1,\ldots,g'_4$ the four displayed new images.
Start with the prescribed values of $V_5$ at all
$1344$ original-sector points. Propagate the equations
\[
                         V_5(g_i u)=g'_iV_5(u)
\]
along the four labeled generator edges. The original
action is transitive on the whole orbit: the first
three generators are transitive on the original
sector, and $r_5$ cycles all five sectors. Hence a
consistent propagation determines every value and
is unique. Reject a conflict or a noninjective
assignment immediately.

For all eighteen pairs the propagation is consistent,
covers the entire $6720$-point set, and yields a
permutation. The certificate then verifies all four
covariance equations pointwise. This proves the
assertion by an explicit finite construction. In
particular these four-generator images generate a
permutation group conjugate to the original image.
No finite-kernel stability was assumed to infer this
conclusion; the intertwining permutations were
constructed and checked directly.

Finally put $N=S^{60}$. By F39 and P29 the complete
$F$ has a primitive in $N$, so $F\in N$ because
$r_5$ normalizes $N$. The complete $R$ is a conjugate
of $F$, and also belongs to $N$. Hence the four old
and new generator images have the same sixty-block
permutation and the same local $L/S$ transitions in
F39's coordinates. The graph construction preserves
each block: this is true on the prescribed original
sector, and each propagation step preserves this
property because $g_i$ and $g'_i$ have the same block
permutation. Transitivity therefore proves that
$V_5$ fixes all sixty blocks, without assuming
beforehand that an arbitrary normalizer preserves
this block system.
Its old-sector components already lie in $S$ by F32.
The covariance equation propagates membership in $S$
from one block to each adjacent block, since the old
and new local transitions have the same class modulo
$S$. Transitivity reaches every block, proving
$V_5\in S^{60}=N\subseteq P'$. Thus the finite
automorphisms here are implemented even by derived-pure
inner conjugation.
\end{proof}

\paragraph{Scope and replay.}
The complete certificate is
\path{agents/pentagon/hurwitz_full_spherical_covariance.py},
with its JSON and NPZ of the eighteen full permutations.
This result concerns the exact prescribed finite
values and does not produce a profinite associator.
The pure order-five primitives are addressed separately
by P29 and F39; the actual-source common-word membership
is the separately proved H9/F38 input.
